\scanpage{23}
\mspage{10}
% archive sheet number: 22
Le sous-groupe $\unsure{\mathscr E_\tau|\mathcal E_\tau}$ de $\mathscr E$ image réciproque de $\mathbb Z/2$ engendré par $\tau$, est le groupe cartographique ``\unsure{triangulé} \unsure{pondéré}'' de Malgoire, ayant comme gén $\tau_0,\tau_1,\tau_\infty$ \deleted{\unsure{relati}} avec
\[
\left\{\begin{aligned}
\deleted{\tau_1\tau_\infty=}{ }l_0&=\tau_\infty\tau_1\\
\deleted{\tau_\infty\tau_0=}{ }l_1&=\tau_0\tau_\infty\\
l_\infty&=\tau_1\tau_0
\end{aligned}\right.\tag*{(17)}
\]
\margin{\uline{NB} \unsure{---}}
$\mathscr E_\tau$ peut être considéré comme le groupe ayant comme \inserted{gén. les $\tau_i$, avec les} seules relations (16) --- on a un hom. signature évident $\tau_i\to-1$ de $\mathscr E_\tau$ dans $\pm1$, dont le noyau est $\pi$ engendré par les $l_i$ décrits par (17)\unsure{....}

\begin{center}---\end{center}
% slanted
\margin{On a les relations}
\margin{$\unsure{\rho}$}
Revenons à une description plus symétrique des autom. extérieurs à lacets de $\mathscr E$. On désignera par $\widehat{\mathscr E}$ le complété profini de $\mathscr E$, extension de $\widehat\Gamma_{0,3}=\Gamma_{0,3}=\unsure{\mathfrak S_3|\mathscr S_3|G_3}\times\mathbb Z/2$ par $\widehat\pi$, et par $\widehat{\widehat{\mathscr E}}$ le groupe des autom. \deleted{extérieurs} à lacets de $\widehat\pi$, extension de $\widehat{\widehat\Gamma}=\operatorname{Aut}\operatorname{ext}_{\unsure{lac|lc}}(\widehat\pi)$ par $\widehat\pi$
\diagram{\(\text{(18)}\quad\)
\begin{tikzcd}[ampersand replacement=\&,column sep=large]
1\arrow[r]\&\widehat\pi\arrow[r]\&\widehat{\widehat{\mathscr E}}\arrow[r]\&\widehat{\widehat\Gamma}\arrow[r]\arrow[d,equal]\&1\\
\&\&\&\operatorname{Aut}\operatorname{ext}_{lac}(\widehat\pi)\&
\end{tikzcd}}
